\chapter{Build: an Execution Algorithm}\label{ch:mx:build-an-execution-algorithm}

A trader hands the algorithm an order to buy 200,000 shares by three o'clock with a limit price. Everything this book has built, the schedule, the
placement, the routing, the measurement, has to run at once, and keep running when the market does something the design did not expect: a halt, a
spike, a kill from the desk. This chapter assembles \texttt{firm.execalgo}, an \omterm{def:mx:benchmark-algorithms:is}{implementation-shortfall algorithm}, from chapters 14, 17, 18 and 25, runs
200 \omterm{def:mx:the-almgren-chriss-framework:parent}{parent orders} through chapter~27's agent market against a TWAP baseline on the same order flow, and measures both with chapter~19's toolkit. The
answer to the question every desk asks, ``how much does it beat TWAP by?'', turns out to depend on the order.

\section{The algorithm as a state machine}

An \omterm{def:mx:benchmark-algorithms:algo}{execution algorithm} (chapter~16) is a program that holds a \omterm{def:mx:the-almgren-chriss-framework:parent}{parent order} and decides, again and again, what \omterm{def:mx:the-almgren-chriss-framework:parent}{child orders} to have in the market.
Written as a state machine it has few states and many events (\cref{fig:mx:build-an-execution-algorithm:states}): a new order starts \emph{working} at
its start time and records the arrival mid; it is \emph{paused} while its stock is halted or paused and resumes when trading does; it ends \emph{done}
when filled, \emph{expired} when its deadline passes with shares left, or \emph{killed} when the desk or a risk control stops it. Every transition and
every refusal is written to an audit log with its time, the state and the reason: the log is how the desk, compliance and the developer reconstruct what
the algorithm did and why.

\begin{omfigure}
\begin{tikzpicture}[font=\footnotesize,st/.style={draw,rounded corners=3pt,minimum width=1.7cm,minimum height=0.65cm,fill=omDef!12},
  end/.style={st,fill=gray!15},>=stealth]
\node[st] (n) at (-0.6,0) {new};
\node[st] (w) at (3.2,0) {working};
\node[st] (p) at (3.2,-2.2) {paused};
\node[end] (d) at (8.4,1.3) {done};
\node[end] (e) at (8.4,0) {expired};
\node[end] (k) at (8.4,-1.3) {killed};
\draw[->,thick] (n) -- node[above] {start time} (w);
\draw[->,thick] ([xshift=-4mm]w.south) -- node[left] {halt, pause} ([xshift=-4mm]p.north);
\draw[->,thick] ([xshift=4mm]p.north) -- node[right] {trading} ([xshift=4mm]w.south);
\draw[->,thick] (w) -- node[above,sloped] {filled} (d);
\draw[->,thick] (w) -- node[above] {deadline} (e);
\draw[->,thick] (w) -- node[above,sloped,pos=0.6] {kill} (k);
\draw[->,thick] (p) -- node[below,sloped] {kill} (k);
\path[->,thick] (w) edge[loop above,looseness=6] node[above,align=center] {every second: schedule,\\ band, cross, rest} (w);
\end{tikzpicture}

\omcaption{The algorithm's states. A working order checks its schedule every second; a halt or pause cancels its orders and stops the clock; the three
end states are final, and every transition is logged.}\label{fig:mx:build-an-execution-algorithm:states}
\end{omfigure}

\section{Schedule, placement and routing together}

The schedule is chapter~14's: the Almgren--Chriss trajectory of holdings $x(t) = X \sinh(\kappa(T - t))/\sinh(\kappa T)$, with the \omterm{def:mx:the-almgren-chriss-framework:urgency}{urgency} given as
$\kappa T$ (zero for a straight line) through \texttt{firm.acexec}'s scheduler interface (Almgren and Chriss 2001; Kissell's textbook describes the same
family of \omterm{def:mx:benchmark-algorithms:is}{implementation-shortfall algorithms}). Every second the algorithm computes what the schedule says should be done by now and what it says should
be done a little later (20 seconds), and bounds both by the \omterm{def:mx:build-an-execution-algorithm:band}{participation band}.

\begin{definition}[Participation band]\label{def:mx:build-an-execution-algorithm:band}
A \emph{participation band}\index{participation band} is a lower and an upper limit on an \omterm{def:mx:benchmark-algorithms:algo}{execution algorithm}'s share of the market's volume since the
order started: whatever its schedule says, the algorithm trades at least enough to stay above the lower limit and never so much as to go above the upper
one.
\end{definition}

\begin{definition}[I-would price]\label{def:mx:build-an-execution-algorithm:iwould}
An \emph{I-would price}\index{I-would price} is a price, set by the trader, at or better than which an algorithm should trade more than its schedule
says, up to its upper participation limit, because the trader would be glad to complete the order there.
\end{definition}

Placement is chapter~17's lesson made operational: rest passively at the near touch for the schedule a little ahead, and cross the spread only for what
the order is behind by more than a tolerance (20 seconds of schedule), so that the passive order has time to fill before it is overtaken. Crossing goes
through chapter~18's sweep when the stock trades on several venues: the displayed quotes at the far touch, best price first, each leg an
immediate-or-cancel child. A limit price caps every child; an \omterm{def:mx:build-an-execution-algorithm:iwould}{I-would price} switches the algorithm to its band's upper edge while the far touch is at or
better than it.

\omcode{code/firm/execalgo/firm_execalgo.py}{170}{209}{The algorithm's decision every second: done or not, the deadline (cross what is left, or expire), the schedule bounded by the participation band, the I-would price, a cross for what is behind by more than the tolerance, and one passive order for the schedule ahead.}\label{lst:mx:build-an-execution-algorithm:check}

\section{Limits, pauses and exceptions}

Every child passes pre-trade controls before it is sent: no child above a maximum size, no price further than a collar from the mid, never more working
and filled than the order. They are the algorithm's own share of the controls a broker must have in front of the market.

\begin{dated}{2026-09}{Market access controls}\label{dat:mx:build-an-execution-algorithm:access}
The SEC adopted Rule 15c3-5 on 3 November 2010, ending unfiltered market access. A broker-dealer with market access must have controls reasonably
designed to prevent the entry of orders that exceed pre-set credit or capital thresholds in the aggregate for each customer and for itself, and of
erroneous orders, by rejecting orders that exceed price or size parameters order by order or over a short period, or that indicate duplicative
orders (17 CFR 240.15c3-5).
\end{dated}

\omcode{code/firm/execalgo/firm_execalgo.py}{94}{108}{The pre-trade controls every child passes: its size, its price against a collar around the consolidated mid, and the order's total; a refusal is logged with its reason.}\label{lst:mx:build-an-execution-algorithm:controls}

The exceptions are driven by the feed and the desk. A trading-action message that leaves the \emph{trading} state (chapter~25's halts and pauses)
cancels every working child and stops the schedule's clock; when trading resumes, the schedule is shifted by the pause, so the order neither rushes to
catch up nor loses its deadline. A kill, sent by the desk or by a risk monitor through the simulator's scheduled calls, cancels everything and ends the
order. Three runs of the same buy of 16,000 shares show them (\cref{fig:mx:build-an-execution-algorithm:stress}, \cref{tab:mx:build-an-execution-algorithm:log}).

\begin{omfigure}
\begin{tikzpicture}
\begin{axis}[width=10.6cm,height=5.8cm,xlabel={\small seconds from the order's start},ylabel={\small filled (thousand shares)},
  tick label style={font=\footnotesize},grid=major,grid style={gray!25},xmin=0,xmax=720,ymin=0,ymax=17,
  legend columns=3,legend style={font=\footnotesize,at={(0.5,-0.28)},anchor=north,draw=none,fill=none,
  /tikz/every even column/.append style={column sep=6pt}},table/col sep=comma]
\addplot[gray,very thick] table[x=t,y=base]{figdata/microstructure/28-build-an-execution-algorithm/stress.csv};
\addplot[omDef,very thick] table[x=t,y=halt]{figdata/microstructure/28-build-an-execution-algorithm/stress.csv};
\addplot[omProp,very thick,dashed] table[x=t,y=spike]{figdata/microstructure/28-build-an-execution-algorithm/stress.csv};
\addplot[omThm,very thick] table[x=t,y=kill]{figdata/microstructure/28-build-an-execution-algorithm/stress.csv};
\legend{normal day,60-second halt,spike with a limit,killed at 300 s}
\end{axis}
\end{tikzpicture}

\omcaption{One buy of 16,000 shares over ten minutes at low \omterm{def:mx:the-almgren-chriss-framework:urgency}{urgency}, on the same order flow: a normal day, a halt from 200 to 260 seconds, a price spike
from another buyer with a limit ten ticks above the arrival mid, and a kill at 300 seconds. Data: \texttt{mx\_execalgo.stress}.}\label{fig:mx:build-an-execution-algorithm:stress}
\end{omfigure}

\begin{table}[ht]
\centering
\small
\begin{tabular}{@{}lrll@{}}
\toprule
run & time (s) & event & detail \\
\midrule
halt & 200.0 & pause & trading state H \\
 & 260.0 & resume & schedule shifted by 60 s \\
 & 661.0 & finish & crossing the last 600 \\
 & 662.0 & done & 16,000 filled \\
spike & 252.0 & limit & far touch 100.18 beyond the limit 100.105 \\
 & 465.0 & limit off & far touch 100.07 \\
 & 587.0 & done & 16,000 filled \\
kill & 299.6 & kill & 8,500 of 16,000 filled \\
\bottomrule
\end{tabular}
\caption{Extracts of the audit log of the three stressed runs (times from the order's start). After the kill, 200 more shares filled: a child
already in flight when the kill was sent. Data: \texttt{mx\_execalgo.stress}.}\label{tab:mx:build-an-execution-algorithm:log}
\end{table}

The halted order finished a minute late, as designed, and paid 2.62 basis points against 2.31 on the normal day. In the spike, another buyer's 60,000
shares took the far touch past the limit within twelve seconds; the algorithm stopped crossing for 213 seconds and kept its passive order at the limit,
and paid 3.79 basis points, against 4.31 for the same order without a limit. The killed order stopped at 8,500 shares and ended at 8,700: a kill
cancels what is resting, not what has already been sent, which is why the desk's position after a kill is read from the fills, never from the order.

\section{Measuring it}

The measure is chapter~19's: the implementation shortfall (Perold 1988) against the arrival mid, all-in (fees and rebates included), attributed by
\texttt{firm.tca} into spread, impact (the order's own push, from a copy of the session without it), timing (the market's own move) and fees. The
experiment is an A/B test with common random numbers: 200 \omterm{def:mx:the-almgren-chriss-framework:parent}{parent orders}, four sizes (4,000 to 32,000 shares over ten minutes, 2.4\% to 15.9\% of the
market's volume) by two urgencies ($\kappa T$ of 0.5 and 3) by 25 seeds, each run with the full algorithm and with a TWAP baseline of the same size on the
same seed (a straight line, a \omterm{def:mx:the-limit-order-book:orders}{market order} every ten seconds for what is due).

\begin{omfigure}
\begin{tikzpicture}
\begin{axis}[width=9.6cm,height=5.6cm,xlabel={\small order size (shares)},ylabel={\small saving against TWAP (bp)},
  xtick={0,1,2,3},xticklabels={4\,000,8\,000,16\,000,32\,000},tick label style={font=\footnotesize},grid=major,grid style={gray!25},
  xmin=-0.4,xmax=3.4,
  legend columns=2,legend style={font=\footnotesize,at={(0.5,-0.3)},anchor=north,draw=none,fill=none,
  /tikz/every even column/.append style={column sep=8pt}},table/col sep=comma]
\addplot[omDef,very thick,mark=*,error bars/.cd,y dir=both,y explicit] table[x=k,y=u0,y error=se0]{figdata/microstructure/28-build-an-execution-algorithm/saved.csv};
\addplot[omThm,very thick,mark=square*,error bars/.cd,y dir=both,y explicit] table[x expr=\thisrow{k}+0.08,y=u1,y error=se1]{figdata/microstructure/28-build-an-execution-algorithm/saved.csv};
\addplot[gray,dashed,domain=-0.4:3.4] {0};
\legend{low urgency ($\kappa T = 0.5$),high urgency ($\kappa T = 3$)}
\end{axis}
\end{tikzpicture}

\omcaption{The all-in saving of the \omterm{def:mx:benchmark-algorithms:is}{implementation-shortfall algorithm} against TWAP on the same order flow, by order size and \omterm{def:mx:the-almgren-chriss-framework:urgency}{urgency}: the mean of 25
paired orders with its standard error. Data: \texttt{mx\_execalgo.ab\_study}.}\label{fig:mx:build-an-execution-algorithm:saved}
\end{omfigure}

Over the 200 orders the algorithm saves $-0.15$ basis points against TWAP, with a 95\% interval of $[-0.44, 0.14]$ clustered by seed: no saving on
average. The average hides a slope (\cref{fig:mx:build-an-execution-algorithm:saved}). A regression of the saving on the size's doubling and a
\omterm{def:mx:the-almgren-chriss-framework:urgency}{high-urgency} dummy gives $1.25 - 0.70\,\log_2(Q/4000) - 0.69\,h$, where $h$ is 1 at high \omterm{def:mx:the-almgren-chriss-framework:urgency}{urgency} (standard errors 0.08, 0.04 and 0.22): the algorithm saves
1.06 basis points on the smallest orders at low \omterm{def:mx:the-almgren-chriss-framework:urgency}{urgency} and loses 1.96 on the largest at high \omterm{def:mx:the-almgren-chriss-framework:urgency}{urgency}, and breaks even near 13,700 shares, about 8\% of
the volume, at low \omterm{def:mx:the-almgren-chriss-framework:urgency}{urgency}.

\begin{omfigure}
\begin{tikzpicture}
\begin{axis}[width=9.6cm,height=5.4cm,ybar,bar width=12pt,ylabel={\small ticks per share},
  xtick={0,1,2,3},xticklabels={spread,impact,timing,fees},tick label style={font=\footnotesize},
  grid=major,grid style={gray!25},enlarge x limits=0.15,
  legend columns=2,legend style={font=\footnotesize,at={(0.5,-0.22)},anchor=north,draw=none,fill=none,
  /tikz/every even column/.append style={column sep=8pt}},table/col sep=comma]
\addplot[fill=omDef!55,draw=omDef] table[x=k,y=is]{figdata/microstructure/28-build-an-execution-algorithm/attribution.csv};
\addplot[fill=gray!45,draw=gray] table[x=k,y=twap]{figdata/microstructure/28-build-an-execution-algorithm/attribution.csv};
\legend{implementation shortfall,TWAP}
\end{axis}
\end{tikzpicture}

\omcaption{\texttt{firm.tca}'s attribution of the two arms' costs, averaged over the 200 pairs (a tick is about a basis point at USD~100). Data:
\texttt{mx\_execalgo.ab\_study}.}\label{fig:mx:build-an-execution-algorithm:attr}
\end{omfigure}

The attribution says why (\cref{fig:mx:build-an-execution-algorithm:attr}). Resting passively earns the algorithm what TWAP pays: $-0.15$ ticks of
spread against $+0.51$, and rebates against fees, $-0.02$ against $+0.30$. It pays it back in impact, 1.73 ticks against 0.93: a resting buy absorbs the
sell orders that would have pushed the price down, the catch-up crosses come in bursts, and \omterm{def:mx:the-almgren-chriss-framework:urgency}{urgency} front-loads the order into the market's memory.
The spread and the fees scale with the shares, the impact with the order's size too, so the passive design wins on small orders and loses on large
ones. This is not a failure of the build: it is the measurement that tells the desk which orders to give it, and which parameters to change for the
others (a lower \omterm{def:mx:the-almgren-chriss-framework:urgency}{urgency}, a smaller tolerance, a wider band).

\section{Rolling it out}

A new algorithm reaches clients the way any model does (Book~7, chapter~21): first in simulation (this chapter), then on a canary, a small share of
the desk's own flow, with a kill criterion written in advance, then in a randomised A/B test against the incumbent, and only then as the default. The
test's size follows from the paired dispersion: the difference between the arms has a standard deviation of 1.33 basis points per order even with
common random numbers, so detecting an average saving of 0.15 with 80\% power at the 5\% level needs $(1.96 + 0.84)^2 \times 1.33^2 / 0.15^2
\approx 620$ orders, and live orders, which do not share their market with a counterfactual, need many more. A per-order effect that depends on size
is easier: stratify by size and test within the stratum where the algorithm is meant to win.

\section{Tutorial: beat TWAP}\label{tut:mx:build-an-execution-algorithm:tut}

\begin{tutorial}
\textbf{Goal.} Assemble an \omterm{def:mx:benchmark-algorithms:is}{implementation-shortfall algorithm}, run it through chapter~27's market against TWAP, measure both with \texttt{firm.tca},
and stress it. \textbf{End state:} \cref{tab:mx:build-an-execution-algorithm:log},
\cref{fig:mx:build-an-execution-algorithm:stress,fig:mx:build-an-execution-algorithm:saved,fig:mx:build-an-execution-algorithm:attr}.
\begin{tutsteps}
\item \textbf{Build.} \texttt{firm\_execalgo.Params(side, qty, start\_s, horizon\_s, urgency, limit, i\_would, band, look\_s, tol\_s)};
\texttt{ISAlgo(params, venues)}; its \texttt{.log}, \texttt{.fills}, \texttt{.state}.
\item \textbf{Run.} \texttt{mx\_execalgo.simulate(seed, agents, controls, calls)}, \texttt{run(params, seed)}, the arms \texttt{IS(qty, urgency)} and
\texttt{TWAP(qty)}.
\item \textbf{Measure.} \texttt{ab\_study()}: the paired savings, \texttt{firm\_tca.cluster\_ols}, the attribution; \texttt{stress()}; draw with
\texttt{fig\_execalgo.py}.
\end{tutsteps}
\textbf{What to change next.} Replace the tolerance rule with chapter~17's dynamic-programming placement; let the \omterm{def:mx:the-almgren-chriss-framework:urgency}{urgency} depend on size; route
passive orders across venues with chapter~18's allocation; run the A/B test on chapter~16's volume curves.
\end{tutorial}

\section{Build: execalgo}\label{bld:mx:build-an-execution-algorithm:execalgo}

\begin{build}
\bfield{Purpose} The firm's \omterm{def:mx:benchmark-algorithms:is}{implementation-shortfall algorithm} on \texttt{firm.exchsim}, for this chapter and for the execution work of Books 11 and 13.
\bfield{Interface} \texttt{Params} (the fields of the tutorial), \texttt{ISAlgo(params, name, venues)} (a \texttt{firm.exchsim} Agent; \texttt{.kill(t\_ns)} for
\texttt{Simulator.schedule\_call}), \texttt{shortfall\_bp(side, fills, arrival, qty, last\_mid)}.
\bfield{Rules} One decision a second; every child through the pre-trade controls; halts cancel and stop the clock; the end states are final; every
transition and refusal in the audit log.
\bfield{Acceptance tests} \texttt{code/firm/execalgo/tests/}: the order completes and \omterm{def:mx:the-almgren-chriss-framework:urgency}{urgency} front-loads it; TWAP only takes; a halt pauses and shifts
the schedule and nothing trades while halted; no fill beyond the limit; the collar refuses and logs; a kill ends the order; the \omterm{def:mx:build-an-execution-algorithm:iwould}{I-would price} trades
ahead of schedule; a cross sweeps two venues' displayed quotes.
\bfield{Stretch} Chapter~17's dynamic-programming placement; dark venues first; a volume-curve schedule; a real-time risk monitor that sends the kill.
\end{build}

\begin{omsources}
\item R.~Almgren and N.~Chriss, ``Optimal execution of portfolio transactions'', \emph{Journal of Risk} 3(2), 2001.
\item A.~F.~Perold, ``The implementation shortfall: paper versus reality'', \emph{Journal of Portfolio Management} 14(3), 1988.
\item R.~Kissell, \emph{The Science of Algorithmic Trading and Portfolio Management}, Academic Press, 2014.
\item SEC Rule 15c3-5, Risk management controls for brokers or dealers with market access, 17 CFR 240.15c3-5 (adopted in Release 34-63241, 2010).
\end{omsources}

\section{Exercises}

\begin{exercise}[$\star$]\label{exo:mx:build-an-execution-algorithm:1}
An order to buy 10,000 shares over 600 seconds with $\kappa T = 0$: how many shares should be done after 300 seconds, and with a tolerance of 20
seconds, from how far behind does the algorithm cross?
\end{exercise}

\begin{exercise}[$\star$]\label{exo:mx:build-an-execution-algorithm:2}
The market has traded 40,000 shares since the order started, of which the algorithm 8,000. With a band of 0 to 30\%, how many more shares may it
trade before trading more?
\end{exercise}

\begin{exercise}[$\star$]\label{exo:mx:build-an-execution-algorithm:3}
From the regression, what is the expected saving of a 16,000-share order at high \omterm{def:mx:the-almgren-chriss-framework:urgency}{urgency}, and at what size does a \omterm{def:mx:the-almgren-chriss-framework:urgency}{low-urgency} order break even?
\end{exercise}

\begin{exercise}[$\star\star$]\label{exo:mx:build-an-execution-algorithm:4}
Why does a passive buy order have impact, when it never crosses the spread?
\end{exercise}

\begin{exercise}[$\star\star$]\label{exo:mx:build-an-execution-algorithm:5}
Why does the algorithm shift its schedule by the length of a halt instead of catching up?
\end{exercise}

\begin{exercise}[$\star\star$]\label{exo:mx:build-an-execution-algorithm:6}
After a kill the order had 200 more shares than when the kill was logged. What should the desk's systems do about it?
\end{exercise}

\begin{exercise}[$\star\star\star$]\label{exo:mx:build-an-execution-algorithm:7}
\emph{Coding.} From the study's pairs, compute the all-in saving before fees are counted, and explain the difference from the all-in figure.
\end{exercise}

\begin{exercise}[$\star\star\star$]\label{exo:mx:build-an-execution-algorithm:8}
\emph{Find the flaw.} ``Our new algorithm beat TWAP by 1 basis point on the 200 orders we gave it last month, so it saves our clients 1 basis point.''
\end{exercise}

\section{Problem: Beat TWAP by How Much?}

\begin{problem}[{Weekend problem --- beat TWAP by how much?}]\label{pb:mx:build-an-execution-algorithm:1}
The head of the execution desk wants to know whether the new \omterm{def:mx:benchmark-algorithms:is}{implementation-shortfall algorithm} should replace TWAP as the default.

\medskip\noindent\textbf{Part I --- The machine.}
\begin{enumerate}
\item Draw the algorithm's states and transitions.
\item What does the audit log record, and for whom?
\item Define a \omterm{def:mx:build-an-execution-algorithm:band}{participation band} and an \omterm{def:mx:build-an-execution-algorithm:iwould}{I-would price}.
\item How are the schedule, the placement and the sweep combined each second?
\item What does the tolerance do, and what happens at the deadline?
\end{enumerate}

\medskip\noindent\textbf{Part II --- Controls and exceptions.}
\begin{enumerate}[resume]
\item What pre-trade controls does every child pass?
\item What does Rule 15c3-5 require of a broker with market access?
\item What does the algorithm do on a halt, and why shift the schedule?
\item What did the limit do in the spike, and what did it save?
\item Why did the killed order end with more shares than it had at the kill?
\end{enumerate}

\medskip\noindent\textbf{Part III --- The experiment.}
\begin{enumerate}[resume]
\item Describe the 200 \omterm{def:mx:the-almgren-chriss-framework:parent}{parent orders} and the TWAP baseline.
\item Why run both arms on the same seeds?
\item What does \texttt{firm.tca} attribute the costs to?
\item \emph{State the named result}: the shortfall saved against TWAP with its confidence interval, and its dependence on \omterm{def:mx:the-almgren-chriss-framework:urgency}{urgency} and order size.
\item Why does the passive design win on small orders and lose on large ones?
\end{enumerate}

\medskip\noindent\textbf{Part IV --- Judgement.}
\begin{enumerate}[resume]
\item How many orders would a live A/B test need, and why more than in simulation?
\item How would you roll the algorithm out?
\item Which orders would you give it, and what would you change for the others?
\item What would you tell the head of the desk?
\item In one sentence: how much does it beat TWAP by?
\end{enumerate}
\end{problem}

\section{Interview questions}

\begin{interviewq}[$\star$ \iqroles{trader}]\label{iq:mx:build-an-execution-algorithm:1}
Your algorithm is 30\% behind schedule with an hour to go and the stock is moving away. What do you do?
\end{interviewq}

\begin{interviewq}[$\star\star$ \iqroles{developer}]\label{iq:mx:build-an-execution-algorithm:2}
How would you structure an \omterm{def:mx:benchmark-algorithms:algo}{execution algorithm}'s code so that it is testable, and what would you log?
\end{interviewq}

\begin{interviewq}[$\star\star$ \iqroles{researcher}]\label{iq:mx:build-an-execution-algorithm:3}
How would you measure whether a new algorithm beats the old one, and how many orders do you need?
\end{interviewq}

\begin{interviewq}[$\star\star$ \iqroles{risk}]\label{iq:mx:build-an-execution-algorithm:4}
Which pre-trade controls belong in the algorithm, and which in a layer the algorithm cannot bypass?
\end{interviewq}

\begin{interviewq}[$\star\star$ \iqroles{bank}]\label{iq:mx:build-an-execution-algorithm:5}
A client asks why your \omterm{def:mx:benchmark-algorithms:is}{implementation-shortfall algorithm} cost more than TWAP on its large orders. What do you answer?
\end{interviewq}

\begin{interviewq}[$\star\star\star$ \iqroles{mle}]\label{iq:mx:build-an-execution-algorithm:6}
You want to learn the algorithm's parameters from its own history of orders. What biases the data, and how do you fix it?
\end{interviewq}
